\section{Desenvolvimento}

O desenvolvimento começou pelo desenho das telas no Figma, com base nos requisitos levantados no questionário. Depois, essas telas foram passadas para um protótipo feito em HTML, CSS e JavaScript, que funciona direto no navegador. Por ser um protótipo, não foi usado servidor nem banco de dados: as contas ficam salvas no navegador e os grupos e tarefas são dados de exemplo.

A aplicação tem dois tipos de usuário, aluno e professor, escolhidos no cadastro. O aluno cria grupos, cadastra tarefas, define o responsável e o prazo de cada uma e atualiza o status conforme o trabalho avança. O professor tem um painel com o andamento dos grupos da turma, onde vê quais grupos têm tarefas atrasadas e pode enviar recados ou lembretes de prazo, que aparecem na página do grupo.

Cada status de tarefa recebeu uma cor (roxo para em andamento, amarelo para pendente, verde para concluída e vermelho para atrasada), e as tarefas atrasadas aparecem em destaque na tela inicial, já que esse recurso foi bem avaliado no questionário. Ficaram de fora nesta versão o armazenamento em servidor e as notificações por e-mail.
